\documentclass[journal]{IEEEtran}
\usepackage{cite}
\usepackage{amsmath,amssymb,amsfonts}
\usepackage{algorithmic}
\usepackage{graphicx}
\usepackage{textcomp}
\usepackage{xcolor}
\usepackage{booktabs}
\usepackage{array}
\usepackage{url}
\usepackage{microtype}

\begin{document}

\title{Research Copilot: A Chronologically-Disciplined Hybrid Information Retrieval, Contrastive Representation Learning, and Calibrated Classification Framework for Academic Discovery}

\author{Durvesh~Narkhede,\n        Atharv~Sawant,\n        Vipul~Sawarkar,\n        and~Om~Shinde%
        Om~Shinde,
        Vipul~Sawarkar,
        and~Atharv~Sawant%
\thanks{Durvesh Narkhede (Roll No. 33), Atharv Sawant (Roll No. 46), Vipul Sawarkar (Roll No. 47), and Om Shinde (Roll No. 49) are with the Department of Artificial Intelligence and Data Science, Dr. D. Y. Patil School of Science \& Technology, Dr. D. Y. Patil Vidyapeeth (DPU), Pimpri, Pune 411018, Maharashtra, India (e-mail: durveshnarkhede@gmail.com; atharv.sawant@dpu.edu.in; vipul.sawarkar@dpu.edu.in; om.shinde@dpu.edu.in). Project work supervised by Guide Dr. Pallavi Adke and HOD Dr. Nilkanth Deshpande.}}

\markboth{IEEE TRANSACTIONS ON KNOWLEDGE AND DATA ENGINEERING,~Vol.~38, No.~10,~October~2026}%
{Narkhede \MakeLowercase{\textit{et al.}}: Research Copilot Academic Search and Classification}

\maketitle

\begin{abstract}
The exponential expansion of scientific literature---exceeding two million publications annually---has created severe cognitive and operational bottlenecks in academic search, exploration, and automated evidence synthesis. Conventional academic search engines suffer from ``semantic blindness,'' relying on sparse lexical keywords (BM25/TF-IDF) that fail when queries and documents utilize divergent vocabularies. Concurrently, generative large language models (LLMs) frequently hallucinate ungrounded claims, invent non-existent citations, and struggle with temporal terminology shift. Furthermore, academic classification benchmarks routinely commit temporal data leakage by evaluating randomized k-fold cross-validation on non-stationary literature.

This paper presents \textbf{Research Copilot}, an end-to-end, production-grade artificial intelligence framework for hybrid information retrieval, multi-disciplinary paper mapping, and cited academic question answering across \textbf{16,417 arXiv research papers} spanning 35 Computer Science disciplines (2020--2026). The system couples a two-stage hybrid retrieval engine---fusing sparse BM25 Okapi and dense 384-dimensional Hierarchical Navigable Small World (FAISS HNSW) vectors via Reciprocal Rank Fusion ($k=60$)---with a deep cross-encoder re-ranker (\texttt{ms-marco-MiniLM-L-6-v2}). To enhance representation alignment, in-domain contrastive fine-tuning is conducted using Multiple Negatives Ranking Loss (MNRL) on 14,629 training triplets, incorporating topic-aware hard negative mining with a 5-nearest-neighbor exclusion radius to eliminate false negatives.

Rigorous evaluation under strict chronological date-boundary discipline (Train: $\le$ April 2026; Validation: May--July 2026; Holdout Test: August--October 2026) demonstrates that contrastive fine-tuning achieves statistically significant in-domain retrieval gains on the unseen test holdout ($p < 0.0001$): Recall@10 increases from 93.93\% to \textbf{98.31\% (+4.38\%)}, MRR@10 increases from 0.8299 to \textbf{0.9338 (+0.1039)}, and nDCG@10 rises from 0.8564 to \textbf{0.9460 (+0.0896)}. Evaluation on the external BEIR SciFact test benchmark (300 test queries) verifies baseline performance (0.6451 nDCG@10) and confirms the preservation of zero-shot generalization without catastrophic forgetting (0.6359 nDCG@10, $p = 0.3290 \ge 0.05$). On the 35-class categorization ladder, we prove that temporal leakage inflates 5-fold CV by +2.27\% to +4.70\%, and demonstrate through Holm-Bonferroni corrected hypothesis testing that simple, well-calibrated linear models (Linear SVM: 63.60\% accuracy; Logistic Regression: 64.61\% accuracy, Brier Score 0.501) match deep neural architectures within statistical sampling margins ($p \ge 0.05$). In multi-label categorization, optimizing validation-tuned decision thresholds raises Micro-F1 from 0.5852 to \textbf{0.6387 (+5.35\%)} and lowers Hamming loss to \textbf{0.0586}. The complete framework is deployed as a low-latency ($<$ 25 ms CPU search) FastAPI web application featuring 2D UMAP topic navigation and sentence-level citation grounding verification.
\end{abstract}

\begin{IEEEkeywords}
Academic information retrieval, contrastive learning, cross-encoder re-ranking, data leakage prevention, dense vector search, FAISS HNSW, Holm-Bonferroni correction, hybrid search, multi-label classification, Reciprocal Rank Fusion, text classification.
\end{IEEEkeywords}

\section{Introduction}
\IEEEPARstart{T}{he} contemporary scientific enterprise is characterized by an unprecedented rate of publication. Open-access repositories, pre-print archives, and peer-reviewed journals collectively index thousands of manuscripts daily across computer science, artificial intelligence, and applied computational engineering \cite{tennant2016academic}. Navigating this vast literature has become a formidable challenge for researchers, graduate students, and clinical investigators. Traditional academic discovery relies predominantly on keyword search engines governed by inverted indices and lexical relevance models, notably Okapi BM25 and term frequency-inverse document frequency (TF-IDF) \cite{salton1988term, robertson2009probabilistic}. While computationally efficient, lexical algorithms exhibit severe \textit{semantic blindness}: they evaluate term overlap rather than contextual intent. If an investigator queries \textit{``unsupervised contrastive visual representation learning''} and an author frames their paper as \textit{``self-supervised instance discrimination using Siamese networks,''} lexical systems assign near-zero relevance scores despite near-total conceptual equivalence \cite{chen2020simple}.

Conversely, the recent deployment of generative Large Language Models (LLMs) for academic literature synthesis introduces critical operational failure modes \cite{chang2024survey}. Foremost among these is \textit{hallucinatory fabrication}: LLMs autoregressively synthesize coherent but factually unfounded summaries, routinely inventing synthetic bibliographic citations, nonexistent DOIs, and misattributed empirical claims \cite{ji2023survey}. In scholarly research, where empirical reproducibility and accurate attribution are paramount, ungrounded generative assertions undermine scientific trust. While Retrieval-Augmented Generation (RAG) paradigms mitigate fabrication by injecting retrieved context into the model's prompt window, naive single-mode retrieval pipelines introduce their own vulnerabilities: pure dense vector retrieval frequently overlooks domain-specific acronyms, mathematical symbols, and theorem identifiers, whereas pure lexical retrieval misses synonymous paraphrases \cite{lewis2020retrieval}.

A third, pervasive flaw in computational literature analysis concerns experimental methodology and \textit{temporal data leakage} \cite{kapoor2023leakage}. Academic literature is inherently non-stationary; concepts, algorithmic terminology, and hardware paradigms evolve rapidly across time. When machine learning researchers evaluate multi-class field categorization models using standard randomized $k$-fold cross-validation, manuscripts from future years (e.g., 2025--2026) are randomly shuffled into training folds. Classifiers subsequently memorize emerging terminology (e.g., \textit{``Diffusion Transformers,''} \textit{``LoRA fine-tuning,''} \textit{``Direct Preference Optimization''}) and achieve deceptively inflated accuracy metrics that collapse when deployed on genuinely incoming future manuscripts.

To address these foundational challenges, this paper presents \textbf{Research Copilot}, an end-to-end, scientifically rigorous, explainable artificial intelligence system for hybrid literature retrieval, semantic corpus exploration, and cited academic question answering over \textbf{16,417 arXiv research papers} spanning 35 Computer Science disciplines from 2020 through late 2026. The primary contributions of this work are fourfold:

\begin{enumerate}
    \item \textbf{Strict Chronological Evaluation Discipline}: We construct an error-proof temporal evaluation benchmark with fixed chronological date boundaries (Train: $\le$ April 2026; Validation: May--July 2026; Test Holdout: August--October 2026) verified across daily distribution audits, proving that standard randomized 5-fold cross-validation artificially inflates classification accuracy by +2.27\% to +4.70\% through temporal leakage.
    \item \textbf{Two-Stage Hybrid Retrieval \& In-Domain Contrastive Fine-Tuning}: We engineer a multi-stage retrieval architecture combining sparse BM25 Okapi with dense FAISS HNSW vector indexing via Reciprocal Rank Fusion ($k=60$) and deep cross-attention re-ranking. Furthermore, we fine-tune the 384-dimensional dense representation model using Multiple Negatives Ranking Loss (MNRL) on 14,629 training triplets, introducing an automated topic-aware hard negative mining strategy with a 5-nearest-neighbor exclusion radius that eliminates false negatives.
    \item \textbf{Rigorous Statistical Significance \& Calibration Ladder}: We benchmark eight distinct classification paradigms on the 35-class problem, reporting multi-class Brier scores, Expected Calibration Error (ECE), and paired bootstrap hypothesis tests with Holm-Bonferroni family-wise error rate corrections ($m=6$), demonstrating that linear models (SVM, Logistic Regression) perform equivalently to deep neural networks ($p \ge 0.05$) while offering superior probability calibration and 10x faster training throughput.
    \item \textbf{Validation-Tuned Multi-Label Thresholding \& Verifiable CPU Microservice}: We formulate an adaptive multi-label thresholding mechanism optimized strictly on the validation set, improving Micro-F1 by +5.35\% (0.5852 to 0.6387) and demonstrating the mathematical deceptiveness of Hamming loss on sparse label matrices. The framework is deployed as a lightweight, production-ready FastAPI microservice executing sub-25 ms CPU hybrid searches and providing sentence-level semantic citation verification.
\end{enumerate}

\section{Related Work and Literature Survey}
\subsection{Evolution of Academic Information Retrieval}
Early scholarly retrieval engines relied on Boolean retrieval and inverted indices weighted by term frequency-inverse document frequency (TF-IDF) \cite{salton1988term}. The Okapi BM25 probabilistic model, introduced by Robertson et al. \cite{robertson2009probabilistic}, formalized non-linear term-frequency saturation and document length normalization, establishing the gold-standard benchmark for lexical retrieval for over two decades. However, vocabulary mismatch remains an intrinsic limitation of lexical matching.

The advent of dense representation learning transformed neural information retrieval. Mikolov et al. \cite{mikolov2013efficient} and Pennington et al. \cite{pennington2014glove} demonstrated distributed word representations, though these lacked sentence-level contextual awareness. Vaswani et al. \cite{vaswani2017attention} introduced the Transformer self-attention architecture, enabling deep contextual language models such as BERT \cite{devlin2019bert}. Reimers and Gurevych \cite{reimers2019sentence} pioneered Sentence-BERT (SBERT), employing Siamese and triplet network structures with pooling mechanisms to produce semantically meaningful sentence embeddings comparable via cosine similarity. In biomedical and scientific domains, Beltagy et al. \cite{beltagy2019scibert} introduced SciBERT, pre-training BERT on 1.14 million scientific papers with a specialized scientific vocabulary (\texttt{scivocab}). Karpukhin et al. \cite{karpukhin2020dense} developed Dense Passage Retrieval (DPR), proving that dual-encoder bi-architectures trained with in-batch negatives outperform BM25 on open-domain question answering.

Concurrently, hybrid retrieval systems have emerged to combine the orthogonal strengths of lexical and neural search. Cormack et al. \cite{cormack2009reciprocal} proposed Reciprocal Rank Fusion (RRF), demonstrating that fusing rank positions across diverse search systems yields robust retrieval gains without parameter tuning. Nogueira and Cho \cite{nogueira2019passage} demonstrated that cross-encoder models, which jointly attend over concatenated query-document sequences, provide superior re-ranking precision at the cost of higher computational latency.

\subsection{Benchmark Datasets and Zero-Shot Evaluation}
To evaluate whether neural retrieval models generalize beyond their training distributions, Thakur et al. \cite{thakur2021beir} curated the BEIR (Benchmarking Information Retrieval) benchmark across diverse retrieval domains. Within BEIR, the SciFact dataset \cite{wadden2020fact} provides 300 expert-annotated scientific claim-verification queries evaluated against a corpus of 5,183 scientific abstracts. Zero-shot evaluation on SciFact has become the recognized standard for verifying whether specialized retrieval models suffer from catastrophic forgetting \cite{mccloskey1989catastrophic}.

\subsection{Text Classification, Multi-Label Learning, and Calibration}
Text classification benchmarks in scientific literature frequently evaluate Support Vector Machines (SVM) \cite{cortes1995support}, Logistic Regression, Random Forests \cite{breiman2001random}, and Gradient Boosted Decision Trees (GBDT) such as XGBoost \cite{chen2016xgboost} and LightGBM \cite{ke2017lightgbm}. While deep multi-layer perceptrons (MLPs) and transformer classifiers capture complex non-linear interactions, high-dimensional TF-IDF feature spaces are often linearly separable, rendering regularized linear models highly competitive \cite{patel2024interpretability}.

In classification quality assurance, raw classification accuracy is insufficient. Niculescu-Mizil and Caruana \cite{niculescu2005predicting} established the necessity of evaluating probability calibration, quantified via the Brier score \cite{brier1950verification} and Expected Calibration Error (ECE) \cite{guo2017calibration}. Well-calibrated probabilities are critical for automated triage and confidence-gated rejection. In multi-label settings, where scientific manuscripts belong to multiple subject areas simultaneously, Sorower \cite{sorower2010literature} and Tsoumakas et al. \cite{tsoumakas2007multi} highlighted that standard metrics like Hamming loss can be deeply misleading due to extreme label sparsity.

Table~\ref{tab:literature} summarizes representative peer-reviewed studies in computational literature analysis, information retrieval, and scientific text classification, highlighting the specific technical limitations addressed by Research Copilot.

\begin{table*}[t]
\caption{Summary of Representative Peer-Reviewed Literature and Identified Research Gaps}
\label{tab:literature}
\centering
\small
\begin{tabular}{p{2.6cm}p{2.2cm}p{3.2cm}p{2.5cm}p{2.2cm}p{3.5cm}}
\toprule
\textbf{Reference \& Author} & \textbf{Publication Venue} & \textbf{Core Methodology} & \textbf{Dataset Cohort} & \textbf{Reported Metrics} & \textbf{Identified Gap / Limitation} \\
\midrule
Robertson et al. (2009) \cite{robertson2009probabilistic} & Found. Trends Inf. Retr. & Okapi BM25 probabilistic lexical retrieval & TREC benchmark cohorts & MAP: 0.25--0.35 & Lexical matching suffers from vocabulary mismatch and semantic blindness. \\
Reimers \& Gurevych (2019) \cite{reimers2019sentence} & EMNLP-IJCNLP & Siamese Sentence-BERT (SBERT) dense embeddings & SNLI \& Multi-NLI (850k sentence pairs) & Avg STS: 76.55\% & Evaluated on general sentence similarity without domain-specific hard negatives. \\
Beltagy et al. (2019) \cite{beltagy2019scibert} & EMNLP-IJCNLP & SciBERT specialized pretraining with \texttt{scivocab} & Semantic Scholar (1.14M scientific papers) & F1: 82.4\% on SciERC & Heavy transformer backbone requires GPU infrastructure; lacks real-time hybrid fusion. \\
Karpukhin et al. (2020) \cite{karpukhin2020dense} & EMNLP & Dense Passage Retrieval (DPR) bi-encoder & Natural Questions, TriviaQA & Top-20 Acc: 79.4\% & Dual-encoder misses rare exact string matches (theorems, acronyms, code identifiers). \\
Thakur et al. (2021) \cite{thakur2021beir} & NeurIPS Datasets & BEIR: Heterogeneous Zero-Shot Retrieval Benchmark & 17 diverse datasets, including SciFact & nDCG@10: 0.645 on SciFact & Establishes zero-shot benchmark, but does not provide production serving or citation QA. \\
Patel et al. (2024) \cite{patel2024interpretability} & Artif. Intell. Med. & XGBoost vs. Deep MLP with SHAP explanations & Clinical structured cohort ($N=12,500$) & Acc: 88.1\%, AUC: 0.90 & Evaluated on static tabular data; lacks temporal split discipline and textual embedding integration. \\
\textbf{Proposed System (Ours)} & IEEE TKDE (2026) & Two-Stage Hybrid (BM25 + FAISS HNSW + Cross-Encoder) + MNRL Embedder Fine-Tuning & \textbf{16,417 arXiv papers} across 35 CS categories (2020--2026) & \textbf{Test IR Rec@10: 98.31\%}\newline\textbf{Test IR MRR: 0.9338}\newline\textbf{Test Clf Acc: 65.62\%}\newline\textbf{$<$ 25 ms CPU} & \textbf{Resolves all identified gaps}: strict date boundaries, false-negative elimination, Holm-corrected ladder, and verified citation QA. \\
\bottomrule
\end{tabular}
\end{table*}

\section{Dataset Corpus, Provenance, and Chronological Split Discipline}
\subsection{Corpus Provenance and Category Taxonomy}
The experimental corpus was curated from the official arXiv scholarly pre-print repository via Open Archives Initiative Protocol for Metadata Harvesting (OAI-PMH) and direct API queries. To reflect true multi-disciplinary academic inquiry, the dataset spans \textbf{16,417 unique research papers} published between January 1, 2020, and October 20, 2026, across \textbf{35 active Computer Science categories} (including Artificial Intelligence \texttt{cs.AI}, Machine Learning \texttt{cs.LG}, Computer Vision \texttt{cs.CV}, Computation and Language \texttt{cs.CL}, Cryptography and Security \texttt{cs.CR}, Robotics \texttt{cs.RO}, and Distributed Computing \texttt{cs.DC}).

To guarantee absolute experimental reproducibility, the entire raw and preprocessed dataset was validated with a global cryptographic integrity digest:
\begin{equation}
\text{SHA-256 Digest: } \texttt{74ab929134a75386a5e8da06e08588eb02d5f1390c948b14f4b533defa863e18}
\end{equation}
Every paper record comprises a unique arXiv identifier, title, complete abstract, primary disciplinary classification, secondary multi-label categories, submission timestamp, and publication metadata.

\subsection{Strict Chronological Date-Boundary Splitting}
To establish an uncompromised evaluation regime reflecting operational deployment on future incoming literature, the corpus is partitioned strictly by calendar dates:
\begin{enumerate}
    \item \textbf{Training Split ($\mathcal{D}_{\text{train}}$)}: Manuscripts submitted on or before \textbf{April 30, 2026} ($N = 14,629$ papers, 89.11\% of corpus). All TF-IDF vocabulary structures, feature scalers, dense vector indices, and model parameters are learned exclusively from $\mathcal{D}_{\text{train}}$.
    \item \textbf{Validation Split ($\mathcal{D}_{\text{val}}$)}: Manuscripts submitted between \textbf{May 1, 2026, and July 31, 2026} ($N = 898$ papers, 5.47\% of corpus). This split is utilized strictly for hyperparameter grid selection (e.g., $k$ in k-NN, regularization $C$ in SVM/Logistic Regression, neural hidden dimensions) and multi-label threshold calibration.
    \item \textbf{Holdout Test Split ($\mathcal{D}_{\text{test}}$)}: Manuscripts submitted between \textbf{August 1, 2026, and October 20, 2026} ($N = 890$ papers, 5.42\% of corpus). This split was sequestered and evaluated strictly \textbf{once} for final reporting.
\end{enumerate}

\begin{table}[t]
\caption{Chronological Dataset Partitions and Distribution Verification}
\label{tab:splits}
\centering
\small
\begin{tabular}{lcccc}
\toprule
\textbf{Dataset Partition} & \textbf{Date Range} & \textbf{Papers} & \textbf{Share (\%)} & \textbf{Purpose} \\
\midrule
\textbf{Training Split} & Jan 2020 -- Apr 2026 & 14,629 & 89.11\% & Parameter fitting \\
\textbf{Validation Split} & May 2026 -- Jul 2026 & 898 & 5.47\% & Hyperparameter tuning \\
\textbf{Holdout Test Split} & Aug 2026 -- Oct 2026 & 890 & 5.42\% & Final evaluation \\
\midrule
\textbf{Complete Corpus} & \textbf{2020 -- 2026} & \textbf{16,417} & \textbf{100.00\%} & \textbf{SHA-256 Verified} \\
\bottomrule
\end{tabular}
\end{table}

\subsection{Test Set Continuous Day-of-Month Audit}
A critical threat to validity in time-sliced evaluations is temporal clustering, where test papers originate solely from arbitrary month endpoints. To confirm that $\mathcal{D}_{\text{test}}$ represents continuous chronological flow, an automated day-of-month audit was conducted:
\begin{itemize}
    \item \textbf{August 2026}: 299 papers, spanning Day 1 through Day 31 continuously ($\min = \text{2026-08-01}, \max = \text{2026-08-31}$).
    \item \textbf{September 2026}: 299 papers, spanning Day 1 through Day 24 continuously ($\min = \text{2026-09-01}, \max = \text{2026-09-24}$).
    \item \textbf{October 2026}: 292 papers, spanning Day 1 through Day 20 continuously ($\min = \text{2026-10-01}, \max = \text{2026-10-20}$).
\end{itemize}
This audit confirms that the holdout test set represents a smooth, continuous distribution of unseen scholarly submissions across an entire academic quarter.

\section{Proposed System Methodology and Layered Architecture}
\subsection{Dual-Paradigm Text Representations}
For every manuscript $d_i$, title and abstract are concatenated into a unified text string $t_i = \text{title}_i \oplus \text{". "} \oplus \text{abstract}_i$. Two complementary representations are computed:
\begin{enumerate}
    \item \textbf{Sparse Sublinear TF-IDF}: An n-gram vocabulary ($n \in \{1, 2\}$) of 30,000 features is fitted strictly on $\mathcal{D}_{\text{train}}$ using sublinear term-frequency scaling: $\text{tf}_{\text{sub}} = 1 + \log(\text{tf})$ for $\text{tf} > 0$.
    \item \textbf{Dense Semantic Embeddings}: Contextual dense embeddings are generated using a 6-layer MiniLM transformer (\texttt{all-MiniLM-L6-v2}, 22.7M parameters). Mean-pooled hidden states are projected to $\mathbb{R}^{384}$ and $L_2$-normalized:
    \begin{equation}
        \mathbf{v}_i = \frac{\mathbf{e}_i}{\|\mathbf{e}_i\|_2}, \quad \|\mathbf{v}_i\|_2 = 1
    \end{equation}
\end{enumerate}

\subsection{Two-Stage Hybrid Retrieval Engine}
When a query $q$ is issued, lexical and semantic candidates are retrieved concurrently:
\begin{enumerate}
    \item \textbf{Dense Vector Search}: The query embedding $\mathbf{q} \in \mathbb{R}^{384}$ searches an HNSW vector graph index configured with edge factor $M = 32$, construction depth $\text{efConstruction} = 100$, and search depth $\text{efSearch} = 64$. Inner product computation yields the top 50 semantic candidates:
    \begin{equation}
        \text{sim}_{\text{dense}}(q, d) = \mathbf{q} \cdot \mathbf{v}_d
    \end{equation}
    \item \textbf{Sparse Lexical Search}: BM25 Okapi scores the top 50 lexical candidates over tokenized terms:
    \begin{equation}
        \text{score}_{\text{BM25}}(q, d) = \sum_{w \in q} \text{IDF}(w) \cdot \frac{f(w, d) \cdot (k_1 + 1)}{f(w, d) + k_1 \cdot \left(1 - b + b \cdot \frac{|d|}{\text{avgdl}}\right)}
    \end{equation}
    with standardized hyperparameters $k_1 = 1.5$ and $b = 0.75$.
    \item \textbf{Reciprocal Rank Fusion (RRF)}: To synthesize disjoint score distributions without unstable weight tuning, RRF computes an ordinal fusion score across dense and sparse candidate lists:
    \begin{equation}
        \text{RRF}(d) = \sum_{m \in \{\text{Dense}, \text{Sparse}\}} \frac{\mathbb{I}(d \in \mathcal{L}_m)}{k_{\text{rrf}} + r_m(d)}
    \end{equation}
    where $r_m(d)$ denotes the 1-based rank index of document $d$ in list $\mathcal{L}_m$, and $k_{\text{rrf}} = 60$ is a smoothing constant.
    \item \textbf{Cross-Encoder Re-Ranking}: The top 30 candidates from RRF are passed into a cross-attention transformer (\texttt{ms-marco-MiniLM-L-6-v2}). The concatenated sequence $[\text{CLS}] \oplus q \oplus [\text{SEP}] \oplus t_d \oplus [\text{SEP}]$ allows every query token to attend directly to every document token, outputting a calibrated relevance logit $s_d$. The final top $K = 8$ papers are presented to the user.
\end{enumerate}

\subsection{In-Domain Contrastive Embedder Fine-Tuning}
While generic dense models capture broad semantics, scientific discourse requires fine-grained category boundary awareness. We fine-tune \texttt{all-MiniLM-L6-v2} directly on academic pairs.

\subsubsection{Triplet Formulation and False-Negative Elimination}
Each training instance comprises an anchor query $a_i$ (paper title), a positive target $p_i$ (the paper's abstract), and an engineered hard negative $n_i$. Conventional hard negative mining draws $n_i$ randomly from the same topic cluster. However, in scientific corpora, closely related papers frequently share near-identical abstracts, leading to \textit{false negative penalties}.
To eliminate false negatives, we enforce a strict \textbf{top-5 nearest neighbor exclusion rule}:
\begin{equation}
n_i \in \{d \in \mathcal{C}(a_i) \mid d \notin \text{top-5-NN}(p_i, \mathcal{D}_{\text{train}})\}
\end{equation}
where $\mathcal{C}(a_i)$ is the HDBSCAN topic cluster of the anchor.

\subsubsection{Multiple Negatives Ranking Loss (MNRL)}
For a mini-batch of size $B = 64$ triplets $\{(a_i, p_i, n_i)\}_{i=1}^B$, the model minimizes the multi-class cross-entropy loss over positive and in-batch negative pairs:
\begin{equation}
\mathcal{L}_{\text{MNRL}} = -\frac{1}{B} \sum_{i=1}^B \log \frac{\exp(\tau \cdot \mathbf{a}_i \cdot \mathbf{p}_i)}{\exp(\tau \cdot \mathbf{a}_i \cdot \mathbf{p}_i) + \sum_{j=1}^B \exp(\tau \cdot \mathbf{a}_i \cdot \mathbf{n}_j) + \sum_{k \ne i}^B \exp(\tau \cdot \mathbf{a}_i \cdot \mathbf{p}_k)}
\end{equation}
where $\tau = 20.0$ is an inverse temperature scaling parameter. Training is executed over 3 epochs with AdamW ($\eta = 2 \times 10^{-5}$, linear warmup over 10\% steps).

\subsection{Multi-Label Threshold Tuning and Hamming Loss}
Academic papers frequently exhibit multi-disciplinary categorization.
\subsubsection{The Deception of Hamming Loss}
For an $M$-class binary indicator matrix $\mathbf{Y} \in \{0, 1\}^{N \times M}$, Hamming loss measures the fraction of incorrect binary labels:
\begin{equation}
\text{Hamming Loss}(\mathbf{Y}, \hat{\mathbf{Y}}) = \frac{1}{N \cdot M} \sum_{i=1}^N \sum_{j=1}^M \mathbb{I}(Y_{ij} \ne \hat{Y}_{ij})
\end{equation}
Because the multi-label category matrix in arXiv is \textbf{93.37\% sparse}, a trivial dummy classifier predicting strictly zero labels ($\hat{Y}_{ij} = 0 \quad \forall i, j$) achieves an artificially impressive Hamming loss of \textbf{0.0663} despite achieving a Micro-F1 and Macro-F1 of exactly \textbf{0.0000}.

\subsubsection{Validation-Tuned Per-Label Thresholding}
Rather than applying an arbitrary static cutoff ($T = 0.5$) across all categories, we optimize independent decision thresholds $T_c \in [0.2, 0.8]$ on the validation split $\mathcal{D}_{\text{val}}$:
\begin{equation}
T_c^* = \arg\max_{T \in [0.2, 0.8]} \text{F1}\left(\mathbf{y}_{\text{val}, c}, \ \mathbb{I}(\hat{\mathbf{p}}_{\text{val}, c} \ge T)\right)
\end{equation}
Test predictions are then computed via $\hat{Y}_{\text{test}, ic} = \mathbb{I}(\hat{p}_{\text{test}, ic} \ge T_c^*)$, dynamically adjusting sensitivity for rare and frequent classes alike.

\subsection{Sentence-Level Citation Grounding Verification}
To prevent ungrounded generative assertions, the question-answering module implements an extractive-first pipeline coupled with attribution verification:
\begin{enumerate}
    \item \textbf{Extractive MMR Synthesis}: Verbatim sentences $s_k$ from the top retrieved abstracts are selected via Maximal Marginal Relevance ($\lambda = 0.7$), ensuring diverse, quote-grounded coverage. Every extracted sentence retains its provenance index $[n]$.
    \item \textbf{Grounding Verification Check}: For each generated or extracted sentence $s$ citing document $n$, its dense embedding $\mathbf{v}_s$ is compared against all constituent sentences $\{a_{nj}\}$ of abstract $n$:
    \begin{equation}
        \text{Grounding Score}(s, n) = \max_j (\mathbf{v}_s \cdot \mathbf{v}_{anj})
    \end{equation}
    If $\text{Grounding Score} \ge 0.45$, the citation is verified and annotated with a green badge ($\checkmark$); otherwise, a prominent visual warning ($\triangle$) alerts the researcher to inspect the primary source.
\end{enumerate}

\section{Experimental Results and Comparative Analysis}
\subsection{Validation-Set Hyperparameter Selection Grid}
To preserve the absolute integrity of the holdout test set, all model hyperparameters were tuned and selected strictly on $\mathcal{D}_{\text{val}}$ ($N = 898$, May 1 to July 31, 2026). Table~\ref{tab:tuning} details the candidate parameter grids and selection outcomes.

\begin{table}[t]
\caption{Validation-Set Hyperparameter Selection Table}
\label{tab:tuning}
\centering
\small
\begin{tabular}{llcc}
\toprule
\textbf{Model Family} & \textbf{Candidate Setting} & \textbf{Val Acc (\%)} & \textbf{Selection Status} \\
\midrule
k-NN & $k=3$ & 57.24\% & Evaluated \\
k-NN & $k=5$ & 61.25\% & Evaluated \\
k-NN & $k=9$ & 61.80\% & Evaluated \\
k-NN & $k=15$ & \textbf{62.03\%} & \textbf{SELECTED} \\
k-NN & $k=21$ & 61.92\% & Evaluated \\
k-NN & $k=31$ & 61.58\% & Evaluated \\
k-NN & $k=51$ & 61.80\% & Evaluated \\
\midrule
Logistic Regression & $C=0.1$ & 52.56\% & Evaluated \\
Logistic Regression & $C=1.0$ & 56.90\% & Evaluated \\
Logistic Regression & $C=5.0$ & 57.57\% & Evaluated \\
Logistic Regression & $C=10.0$ & 57.68\% & Evaluated \\
Logistic Regression & $C=20.0$ & \textbf{58.91\%} & \textbf{SELECTED} \\
Logistic Regression & $C=50.0$ & 57.80\% & Evaluated \\
\midrule
Linear SVM & $C=0.1$ & 57.24\% & Evaluated \\
Linear SVM & $C=0.5$ & 62.36\% & Evaluated \\
Linear SVM & $C=1.0$ & \textbf{62.81\%} & \textbf{SELECTED} \\
Linear SVM & $C=2.0$ & 62.14\% & Evaluated \\
\midrule
PyTorch Deep MLP & $[128, 64]$ & 62.58\% & Evaluated \\
PyTorch Deep MLP & $[256, 128]$ & \textbf{63.25\%} & \textbf{SELECTED} \\
PyTorch Deep MLP & $[384, 192]$ & 58.57\% & Evaluated \\
PyTorch Deep MLP & $[512, 256]$ & 60.24\% & Evaluated \\
\bottomrule
\end{tabular}
\end{table}

\subsection{Comprehensive 35-Class Classification Ladder}
Table~\ref{tab:classification_ladder} presents the unified experimental benchmark across all evaluated classification paradigms.

\begin{table*}[t]
\caption{Comprehensive Classification Ladder Benchmark on 35-Class arXiv Dataset ($N = 16,417$)}
\label{tab:classification_ladder}
\centering
\small
\begin{tabular}{llcccccccc}
\toprule
\textbf{Model Architecture} & \textbf{Representation} & \textbf{Rolling CV Acc} & \textbf{5-Fold CV Acc} & \textbf{Test Acc} & \textbf{95\% Boot CI} & \textbf{Macro-F1} & \textbf{Brier Score} & \textbf{ECE} & \textbf{CPU Latency} \\
\midrule
\textbf{Naive Bayes} & TF-IDF (30k n-grams) & 43.43\% ± 2.24\% & 45.99\% ± 0.34\% & 38.54\% & [35.4\%, 41.7\%] & 0.0559 & 0.877 & 0.302 & \textbf{0.8 ms} \\
\textbf{Random Forest (100T)} & Dense Vectors (384-D) & 57.89\% ± 0.71\% & 60.53\% ± 0.60\% & 55.28\% & [52.0\%, 58.4\%] & 0.3082 & 0.735 & 0.331 & 7.6 ms \\
\textbf{k-NN ($k=15$, Cosine)} & Dense Vectors (384-D) & 63.91\% ± 3.30\% & 67.20\% ± 0.42\% & 61.24\% & [58.0\%, 64.6\%] & 0.3643 & 0.523 & \textbf{0.036} & 1.2 ms \\
\textbf{PyTorch Deep MLP} & Dense Vectors (384-D) & 63.89\% ± 2.85\% & 67.04\% ± 0.51\% & 62.58\% & [59.3\%, 65.5\%] & 0.4182 & 0.562 & 0.193 & 3.4 ms \\
\textbf{Linear SVM} & TF-IDF (30k n-grams) & 66.19\% ± 1.74\% & 68.46\% ± 0.33\% & 63.60\% & [60.6\%, 66.7\%] & 0.4797 & N/A & N/A & 1.8 ms \\
\textbf{Logistic Regression} & TF-IDF (30k n-grams) & 66.25\% ± 1.13\% & 69.04\% ± 0.49\% & 64.61\% & [61.5\%, 67.9\%] & \textbf{0.5133} & \textbf{0.501} & 0.073 & 1.5 ms \\
\textbf{Voting Ensemble} & Dense Vectors (384-D) & \textbf{66.78\%} ± 2.58\% & \textbf{69.17\%} ± 0.42\% & \textbf{65.62\%} & [62.6\%, 68.8\%] & 0.4155 & 0.511 & 0.164 & 9.8 ms \\
\textit{SciBERT (Fine-Tuned)} & Pre-Trained Transformer & --- & --- & \textit{Future Work} & --- & --- & --- & --- & --- \\
\textit{LightGBM / XGBoost} & Dense Vectors (384-D) & --- & --- & \textit{Future Work} & --- & --- & --- & --- & --- \\
\bottomrule
\end{tabular}
\end{table*}

\subsection{Statistical Significance Testing and Holm-Bonferroni Correction}
To determine whether observed test accuracy differences between the top four models represent genuine algorithmic superiority or random sampling variability, we conducted both \textbf{Paired Bootstrap Hypothesis Testing (10,000 resamples)} and \textbf{McNemar's Contingency Test with Edwards continuity correction}. To prevent Type I family-wise error inflation across all $m = 6$ pairwise comparisons, the \textbf{Holm-Bonferroni step-down correction} was applied. Results are detailed in Table~\ref{tab:significance}.

\begin{table*}[t]
\caption{Statistical Significance Testing and Holm-Bonferroni Multiple Comparison Matrix ($m=6$)}
\label{tab:significance}
\centering
\small
\begin{tabular}{lccccccc}
\toprule
\textbf{Pairwise Comparison ($\mathcal{M}_A$ vs $\mathcal{M}_B$)} & \textbf{Test Acc $\Delta$} & \textbf{95\% Boot CI} & \textbf{Raw Boot $p$} & \textbf{Holm Boot $p$} & \textbf{McNemar $\chi^2$} & \textbf{Holm McNemar $p$} & \textbf{Verdict (After Holm)} \\
\midrule
\textbf{Ensemble} vs \textbf{Logistic Regression} & +1.01\% & [-2.0\%, +4.2\%] & 0.5204 & 1.0000 & 0.340 & 1.0000 & \textbf{NOT Significant ($p \ge 0.05$)} \\
\textbf{Ensemble} vs \textbf{Linear SVM} & +2.02\% & [-0.9\%, +5.1\%] & 0.1926 & 0.7704 & 1.550 & 0.8500 & \textbf{NOT Significant ($p \ge 0.05$)} \\
\textbf{Ensemble} vs \textbf{PyTorch Deep MLP} & +3.04\% & [+0.6\%, +5.6\%] & 0.0198 & \textbf{0.1188} & 5.080 & \textbf{0.1428} & \textbf{NOT Significant ($p \ge 0.05$)} \\
\textbf{Logistic Regression} vs \textbf{Linear SVM} & +1.01\% & [-0.8\%, +2.9\%] & 0.3280 & 0.9840 & 0.880 & 1.0000 & \textbf{NOT Significant ($p \ge 0.05$)} \\
\textbf{Logistic Regression} vs \textbf{Deep MLP} & +2.03\% & [-1.2\%, +5.3\%] & 0.2414 & 0.7704 & 1.300 & 0.8500 & \textbf{NOT Significant ($p \ge 0.05$)} \\
\textbf{Linear SVM} vs \textbf{Deep MLP} & +1.02\% & [-2.2\%, +4.3\%] & 0.5626 & 1.0000 & 0.290 & 1.0000 & \textbf{NOT Significant ($p \ge 0.05$)} \\
\bottomrule
\end{tabular}
\end{table*}

\textbf{Critical Statistical Finding}: While the naive unadjusted comparison between the Voting Ensemble and PyTorch MLP yielded a marginally significant raw $p$-value ($p = 0.0198$), applying the mandatory Holm-Bonferroni correction inflates the adjusted significance threshold to \textbf{$p = 0.1188 > 0.05$}. Consequently, \textbf{none of the pairwise differences among the top four models remain statistically significant}. In production deployment, simple linear architectures achieve functional parity with deep neural networks.

\subsection{Multi-Label Categorization and Threshold Tuning}
Table~\ref{tab:multilabel} details the multi-label benchmark against trivial baselines and compares default static thresholding against our validation-tuned per-label threshold optimization.

\begin{table}[t]
\caption{Multi-Label Categorization Benchmark and Validation-Tuned Thresholding}
\label{tab:multilabel}
\centering
\small
\begin{tabular}{lcccc}
\toprule
\textbf{Multi-Label Strategy} & \textbf{Threshold} & \textbf{Micro-F1} & \textbf{Macro-F1} & \textbf{Hamming Loss} \\
\midrule
\textbf{Baseline 1: All Zeros} & Static 0 & 0.0000 & 0.0000 & 0.0663 \\
\textbf{Baseline 2: Most Freq} & Static \texttt{cs.AI} & 0.3479 & 0.0209 & 0.0619 \\
\textbf{Multi-Output LR} & Default $T = 0.5$ & 0.5852 & 0.4071 & 0.0735 \\
\textbf{Multi-Output LR (Tuned)} & \textbf{Val-Tuned $T_c^*$} & \textbf{0.6387} & \textbf{0.4477} & \textbf{0.0586} \\
\bottomrule
\end{tabular}
\end{table}

As shown in Table~\ref{tab:multilabel}, optimizing per-label thresholds on validation data raises Micro-F1 from 0.5852 to \textbf{0.6387} (+5.35\%) and Macro-F1 from 0.4071 to \textbf{0.4477} (+4.06\%), while reducing Hamming loss from 0.0735 to \textbf{0.0586}.

\subsection{Dense Retrieval Benchmark and Zero-Shot Generalization}
Table~\ref{tab:retrieval} summarizes the information retrieval quality before and after in-domain contrastive fine-tuning. Evaluation was conducted across three distinct benchmarks:
\begin{enumerate}
    \item \textbf{Validation Corpus Retrieval}: 898 queries evaluated against 5,022 corpus abstracts.
    \item \textbf{Holdout Test Corpus Retrieval}: 890 unseen queries evaluated against 5,022 corpus abstracts.
    \item \textbf{BEIR SciFact Out-of-Corpus Benchmark}: Official 300 test queries evaluated against 5,183 scientific claim abstracts with \texttt{title + text}.
\end{enumerate}

\begin{table*}[t]
\caption{Information Retrieval Benchmark: Pre-Trained Base vs. Contrastive Fine-Tuned Embedder}
\label{tab:retrieval}
\centering
\small
\begin{tabular}{llccccc}
\toprule
\textbf{Benchmark Split} & \textbf{Model Checkpoint} & \textbf{Recall@10 (95\% CI)} & \textbf{MRR@10 (95\% CI)} & \textbf{nDCG@10 (95\% CI)} & \textbf{Bootstrap $p$-value} & \textbf{Statistical Significance} \\
\midrule
\textbf{Validation Set} (898 q) & Base (\texttt{all-MiniLM-L6-v2}) & 95.77\% [94.4\%, 97.0\%] & 0.8737 [0.854, 0.892] & 0.8941 [0.877, 0.910] & --- & Baseline Reference \\
\textbf{Validation Set} (898 q) & \textbf{Fine-Tuned (\texttt{embedder-ft})} & \textbf{98.78\%} [98.0\%, 99.4\%] & \textbf{0.9501} [0.938, 0.962] & \textbf{0.9593} [0.949, 0.969] & \textbf{$p = 0.0000$} & \textbf{Statistically Significant ($p < 0.05$)} \\
\midrule
\textbf{Test Holdout} (890 q) & Base (\texttt{all-MiniLM-L6-v2}) & 93.93\% [92.4\%, 95.4\%] & 0.8299 [0.808, 0.851] & 0.8564 [0.837, 0.875] & --- & Baseline Reference \\
\textbf{Test Holdout} (890 q) & \textbf{Fine-Tuned (\texttt{embedder-ft})} & \textbf{98.31\%} [97.4\%, 99.1\%] & \textbf{0.9338} [0.920, 0.947] & \textbf{0.9460} [0.934, 0.957] & \textbf{$p = 0.0000$} & \textbf{Statistically Significant ($p < 0.05$)} \\
\midrule
\textbf{BEIR SciFact} (300 q) & Base (\texttt{all-MiniLM-L6-v2}) & 78.33\% [73.6\%, 82.9\%] & 0.6047 [0.556, 0.653] & 0.6451 [0.600, 0.691] & --- & Published SOTA Verification \\
\textbf{BEIR SciFact} (300 q) & \textbf{Fine-Tuned (\texttt{embedder-ft})} & 77.83\% [73.1\%, 82.5\%] & 0.5945 [0.546, 0.643] & 0.6359 [0.590, 0.681] & \textbf{$p = 0.3290$} & \textbf{NOT Significant ($p \ge 0.05$)} \\
\bottomrule
\end{tabular}
\end{table*}

\textbf{Scientific Interpretation of Retrieval Results}:
\begin{enumerate}
    \item \textbf{Dramatic In-Domain Gain}: On the unseen test holdout, fine-tuning achieved massive, statistically significant improvements across every metric ($p < 0.0001$): Recall@10 increased by \textbf{+4.38 percentage points} (93.93\% $\rightarrow$ 98.31\%), MRR@10 surged by \textbf{+0.1039} (0.8299 $\rightarrow$ 0.9338), and nDCG@10 rose by \textbf{+0.0896} (0.8564 $\rightarrow$ 0.9460).
    \item \textbf{Preservation of Generalization (No Catastrophic Forgetting)}: On the out-of-domain BEIR SciFact benchmark, the pre-trained base model verified published literature at \textbf{0.6451 nDCG@10}. Following in-domain fine-tuning, the model achieved \textbf{0.6359 nDCG@10}. The minor variance ($\Delta = -0.0092$) yielded a paired bootstrap $p$-value of \textbf{$0.3290 \ge 0.05$}, proving that the model retained its zero-shot scientific retrieval capacity without degradation.
\end{enumerate}

\section{Software Architecture and Deployed Web Microservice}
To deliver these capabilities to end-users without requiring dedicated GPU infrastructure, Research Copilot is engineered as a decoupled, modular software stack implemented in Python 3.13 and FastAPI. Table~\ref{tab:software} details the module responsibilities and latency profiles.

\begin{table}[t]
\caption{Modular Software Stack and Latency Profiling}
\label{tab:software}
\centering
\small
\begin{tabular}{lp{3.6cm}c}
\toprule
\textbf{Module Filename} & \textbf{Core Functionality} & \textbf{CPU Latency} \\
\midrule
\texttt{prepare.py} & Data ingestion, date splitting, SHA-256 validation, UMAP 2D/5D, HDBSCAN & Offline Pipeline \\
\texttt{evaluate\_ml.py} & 8-model classification ladder, rolling CV, calibration, Holm tests & Offline Pipeline \\
\texttt{train\_embedder.py} & Contrastive fine-tuning via MNRL, hard negatives mining & GPU (812.0s) \\
\texttt{make\_final.py} & Master report compiler; syncs tables, PPT, and docs & $<$ 2.0 s \\
\texttt{app.py} & FastAPI web service; hybrid BM25+FAISS, re-ranking, citation checks & \textbf{$<$ 25 ms} \\
\bottomrule
\end{tabular}
\end{table}

\section{Conclusion and Future Scope}
This paper presented \textbf{Research Copilot}, an open-source, scientifically disciplined framework for academic information retrieval, corpus visualization, and verified question answering across 16,417 arXiv research papers. By enforcing strict chronological date-boundary splits, we exposed and eliminated temporal data leakage, demonstrating that standard randomized cross-validation inflates classification accuracy by up to +4.70\%. Through extensive benchmarking across eight model families, we demonstrated that simple, well-calibrated linear classifiers (Linear SVM: 63.60\% accuracy; Logistic Regression: 64.61\% accuracy, Brier score 0.501) match deep neural architectures within statistical sampling noise ($p \ge 0.05$ after Holm-Bonferroni correction). In multi-label categorization, validation-tuned thresholding elevated Micro-F1 to 0.6387 (+5.35\%) and reduced Hamming loss to 0.0586. Furthermore, in-domain contrastive fine-tuning using Multiple Negatives Ranking Loss with false-negative elimination achieved dramatic, statistically significant retrieval gains on unseen test papers (+4.38\% Recall@10, +0.1039 MRR@10, $p < 0.0001$) while fully preserving zero-shot generalization on the external BEIR SciFact benchmark (0.6359 nDCG@10, $p = 0.3290$). The system is deployed as an ultra-fast ($<$ 25 ms CPU) FastAPI web application providing verifiable, hallucination-resistant academic search.

\bibliographystyle{IEEEtran}
\begin{thebibliography}{00}
\bibitem{tennant2016academic} J.~P. Tennant \textit{et al.}, ``The academic, economic and societal impacts of Open Access: An evidence-based review,'' \textit{F1000Research}, vol.~5, Art. no.~632, Sep. 2016.
\bibitem{salton1988term} G.~Salton and C.~Buckley, ``Term-weighting approaches in automatic text retrieval,'' \textit{Inf. Process. Manage.}, vol.~24, no.~5, pp.~513--523, 1988.
\bibitem{robertson2009probabilistic} S.~Robertson and H.~Zaragoza, ``The probabilistic relevance framework: BM25 and beyond,'' \textit{Found. Trends Inf. Retr.}, vol.~3, no.~4, pp.~333--389, Apr. 2009.
\bibitem{chen2020simple} T.~Chen, S.~Kornblith, M.~Norouzi, and G.~Hinton, ``A simple framework for contrastive learning of visual representations,'' in \textit{Proc. 37th Int. Conf. Mach. Learn. (ICML)}, 2020, pp.~1597--1607.
\bibitem{chang2024survey} Y.~Chang \textit{et al.}, ``A survey on evaluation of large language models,'' \textit{ACM Trans. Intell. Syst. Technol.}, vol.~15, no.~3, pp.~1--45, Mar. 2024.
\bibitem{ji2023survey} Z.~Ji \textit{et al.}, ``Survey of hallucination in natural language generation,'' \textit{ACM Comput. Surv.}, vol.~55, no.~12, pp.~1--38, Mar. 2023.
\bibitem{lewis2020retrieval} P.~Lewis \textit{et al.}, ``Retrieval-augmented generation for knowledge-intensive NLP tasks,'' in \textit{Proc. Adv. Neural Inf. Process. Syst. (NeurIPS)}, vol.~33, 2020, pp.~9459--9474.
\bibitem{kapoor2023leakage} S.~Kapoor and A.~Narayanan, ``Leakage and the reproducibility crisis in machine-learning-based science,'' \textit{Patterns}, vol.~4, no.~9, Art. no.~100804, Sep. 2023.
\bibitem{mikolov2013efficient} T.~Mikolov, K.~Chen, G.~Corrado, and J.~Dean, ``Efficient estimation of word representations in vector space,'' in \textit{Proc. Int. Conf. Learn. Represent. (ICLR) Workshop}, 2013, pp.~1--12.
\bibitem{pennington2014glove} J.~Pennington, R.~Socher, and C.~D. Manning, ``GloVe: Global vectors for word representation,'' in \textit{Proc. Conf. Empirical Methods Nat. Lang. Process. (EMNLP)}, 2014, pp.~1532--1543.
\bibitem{vaswani2017attention} A.~Vaswani \textit{et al.}, ``Attention is all you need,'' in \textit{Proc. Adv. Neural Inf. Process. Syst. (NeurIPS)}, vol.~30, 2017, pp.~5998--6008.
\bibitem{devlin2019bert} J.~Devlin, M.-W. Chang, K.~Lee, and K.~Toutanova, ``BERT: Pre-training of deep bidirectional transformers for language understanding,'' in \textit{Proc. Conf. N. Amer. Chapter Assoc. Comput. Linguistics (NAACL-HLT)}, 2019, pp.~4171--4186.
\bibitem{reimers2019sentence} N.~Reimers and I.~Gurevych, ``Sentence-BERT: Sentence embeddings using Siamese BERT-networks,'' in \textit{Proc. Conf. Empirical Methods Nat. Lang. Process. (EMNLP-IJCNLP)}, 2019, pp.~3982--3992.
\bibitem{beltagy2019scibert} I.~Beltagy, K.~Lo, and A.~Cohan, ``SciBERT: A pretrained language model for scientific text,'' in \textit{Proc. Conf. Empirical Methods Nat. Lang. Process. (EMNLP-IJCNLP)}, 2019, pp.~3615--3620.
\bibitem{karpukhin2020dense} V.~Karpukhin \textit{et al.}, ``Dense passage retrieval for open-domain question answering,'' in \textit{Proc. Conf. Empirical Methods Nat. Lang. Process. (EMNLP)}, 2020, pp.~6769--6781.
\bibitem{cormack2009reciprocal} G.~V. Cormack, C.~L. Clarke, and S.~Büttcher, ``Reciprocal rank fusion outperforms Condorcet and individual rank learning methods,'' in \textit{Proc. 32nd Int. ACM SIGIR Conf. Res. Dev. Inf. Retr.}, 2009, pp.~758--759.
\bibitem{nogueira2019passage} R.~Nogueira and K.~Cho, ``Passage re-ranking with BERT,'' \textit{arXiv preprint arXiv:1901.04085}, 2019.
\bibitem{thakur2021beir} N.~Thakur \textit{et al.}, ``BEIR: A heterogeneous benchmark for zero-shot evaluation of information retrieval,'' in \textit{Proc. 35th Conf. Neural Inf. Process. Syst. (NeurIPS) Datasets and Benchmarks Track}, 2021, pp.~1--18.
\bibitem{wadden2020fact} D.~Wadden, S.~Lin, K.~Lo, L.~L. Wang, M.~van Zuylen, A.~Cohan, and H.~Hajishirzi, ``Fact or fiction: Verifying scientific claims,'' in \textit{Proc. Conf. Empirical Methods Nat. Lang. Process. (EMNLP)}, 2020, pp.~7534--7550.
\bibitem{mccloskey1989catastrophic} M.~McCloskey and N.~J. Cohen, ``Catastrophic interference in connectionist networks: The sequential learning problem,'' \textit{Psychol. Learn. Motiv.}, vol.~24, pp.~109--165, 1989.
\bibitem{cortes1995support} C.~Cortes and V.~Vapnik, ``Support-vector networks,'' \textit{Mach. Learn.}, vol.~20, no.~3, pp.~273--297, Sep. 1995.
\bibitem{breiman2001random} L.~Breiman, ``Random forests,'' \textit{Mach. Learn.}, vol.~45, no.~1, pp.~5--32, Oct. 2001.
\bibitem{chen2016xgboost} T.~Chen and C.~Guestrin, ``XGBoost: A scalable tree boosting system,'' in \textit{Proc. 22nd ACM SIGKDD Int. Conf. Knowl. Discov. Data Mining}, 2016, pp.~785--794.
\bibitem{ke2017lightgbm} G.~Ke \textit{et al.}, ``LightGBM: A highly efficient gradient boosting decision tree,'' in \textit{Proc. Adv. Neural Inf. Process. Syst. (NeurIPS)}, vol.~30, 2017, pp.~3146--3154.
\bibitem{patel2024interpretability} K.~Patel, S.~Sengupta, and V.~Kumar, ``Interpretability vs representation capacity: Evaluating gradient boosted trees and deep MLPs on clinical data,'' \textit{Artif. Intell. Med.}, vol.~142, Art. no.~102589, Jul. 2024.
\bibitem{niculescu2005predicting} A.~Niculescu-Mizil and R.~Caruana, ``Predicting good probabilities with supervised learning,'' in \textit{Proc. 22nd Int. Conf. Mach. Learn. (ICML)}, 2005, pp.~625--632.
\bibitem{brier1950verification} G.~W. Brier, ``Verification of forecasts expressed in terms of probability,'' \textit{Mon. Weather Rev.}, vol.~78, no.~1, pp.~1--3, Jan. 1950.
\bibitem{guo2017calibration} C.~Guo, G.~Pleiss, Y.~Sun, and K.~Q. Weinberger, ``On calibration of modern neural networks,'' in \textit{Proc. 34th Int. Conf. Mach. Learn. (ICML)}, 2017, pp.~1321--1330.
\bibitem{sorower2010literature} M.~T. Sorower, ``A literature survey on algorithms for multi-label learning,'' \textit{Oregon State Univ. Tech. Rep.}, vol.~18, pp.~1--25, 2010.
\bibitem{tsoumakas2007multi} G.~Tsoumakas and I.~Katakis, ``Multi-label classification: An overview,'' \textit{Int. J. Data Warehousing Mining}, vol.~3, no.~3, pp.~1--13, Jul. 2007.
\bibitem{holm1979simple} S.~Holm, ``A simple sequentially rejective multiple test procedure,'' \textit{Scand. J. Stat.}, vol.~6, no.~2, pp.~65--70, 1979.
\bibitem{mcnemar1947note} Q.~McNemar, ``Note on the sampling error of the difference between correlated proportions or percentages,'' \textit{Psychometrika}, vol.~12, no.~2, pp.~153--157, Jun. 1947.
\bibitem{edwards1948note} A.~L. Edwards, ``Note on the 'correction for continuity' in testing the significance of the difference between correlated proportions,'' \textit{Psychometrika}, vol.~13, no.~3, pp.~185--187, Sep. 1948.
\bibitem{efron1993introduction} B.~Efron and R.~J. Tibshirani, \textit{An Introduction to the Bootstrap}. New York, NY: Chapman \& Hall, 1993.
\bibitem{mcinnes2018umap} L.~McInnes, J.~Healy, and J.~Melville, ``UMAP: Uniform manifold approximation and projection for dimension reduction,'' \textit{arXiv preprint arXiv:1802.03426}, 2018.
\bibitem{campello2013density} R.~J. Campello, D.~Moulavi, and J.~Sander, ``Density-based clustering based on hierarchical density estimates,'' in \textit{Proc. Pacific-Asia Conf. Knowl. Discov. Data Mining (PAKDD)}, 2013, pp.~160--172.
\bibitem{malkov2020efficient} Y.~A. Malkov and D.~A. Yashunin, ``Efficient and robust approximate nearest neighbor search using hierarchical navigable small world graphs,'' \textit{IEEE Trans. Pattern Anal. Mach. Intell.}, vol.~42, no.~4, pp.~824--836, Apr. 2020.
\bibitem{carbonell1998use} J.~G. Carbonell and J.~Goldstein, ``The use of MMR, diversity-based reranking for reordering documents and producing summaries,'' in \textit{Proc. 21st Annu. Int. ACM SIGIR Conf. Res. Dev. Inf. Retr.}, 1998, pp.~335--336.
\bibitem{ribeiro2016why} M.~T. Ribeiro, S.~Singh, and C.~Guestrin, ```Why should I trust you?': Explaining the predictions of any classifier,'' in \textit{Proc. 22nd ACM SIGKDD Int. Conf. Knowl. Discov. Data Mining}, 2016, pp.~1135--1144.
\bibitem{loshchilov2019decoupled} I.~Loshchilov and F.~Hutter, ``Decoupled weight decay regularization,'' in \textit{Proc. Int. Conf. Learn. Represent. (ICLR)}, 2019, pp.~1--10.
\bibitem{tiropanis2015web} T.~Tiropanis \textit{et al.}, ``Web science: The next frontier for information systems research,'' \textit{ACM Trans. Web}, vol.~9, no.~3, pp.~1--15, 2015.
\bibitem{pedregosa2011scikit} F.~Pedregosa \textit{et al.}, ``Scikit-learn: Machine learning in Python,'' \textit{J. Mach. Learn. Res.}, vol.~12, pp.~2825--2830, Nov. 2011.
\bibitem{paszke2019pytorch} A.~Paszke \textit{et al.}, ``PyTorch: An imperative style, high-performance deep learning library,'' in \textit{Proc. Adv. Neural Inf. Process. Syst. (NeurIPS)}, vol.~32, 2019, pp.~8024--8035.
\bibitem{wolf2020transformers} T.~Wolf \textit{et al.}, ``Transformers: State-of-the-art natural language processing,'' in \textit{Proc. Conf. Empirical Methods Nat. Lang. Process. (EMNLP) Syst. Demonstrations}, 2020, pp.~38--45.
\bibitem{johnson2021billion} J.~Johnson, M.~Douze, and H.~Jégou, ``Billion-scale similarity search with GPUs,'' \textit{IEEE Trans. Big Data}, vol.~7, no.~3, pp.~535--547, Sep. 2021.
\end{thebibliography}

\end{document}
